%*************************************************
% همه بسته‌ها و تنظیمات کلی اینجاست (برگرفته از قالب دانشگاه صنعتی اصفهان)
%*************************************************
\usepackage{amsthm,amssymb,amsmath}      % ریاضیات
\usepackage{graphicx}                     % درج تصویر
\usepackage[a4paper]{geometry}            % حاشیه‌ها
\usepackage{titlesec}                     % ظاهر عنوان‌ها
\usepackage{setspace}                     % فاصله خطوط
\usepackage[stable,bottom]{footmisc}      % پاورقی در انتهای صفحه
\usepackage{zref-perpage}                 % شماره‌گذاری پاورقی در هر صفحه از اول
\zmakeperpage[1]{footnote}

\usepackage{xepersian}

% فونت اصلی فارسی
\settextfont{XB Zar}[
  Path = ./,
  Extension = .ttf,
  BoldFont = XB ZarBd,
  ItalicFont = XB ZarIt,
  BoldItalicFont = XB ZarBdIt
]

% پاورقی‌ها چپ‌به‌راست (برای واژه‌های انگلیسی)
\setfootnoteLR

\makeatletter
\def\@makeLTRfnmark{\hbox{\@textsuperscript{\latinfont\@thefnmark}}}
\renewcommand\@makefntext[1]{%
    \parindent 1em%
    \noindent
    \hb@xt@1.8em{\hss\if@RTL\@makefnmark\else\@makeLTRfnmark\fi}#1}
\makeatother

% جداکننده شماره بخش‌ها: خط تیره  (۱-۲-)
\SepMark{-}

% ظاهر عنوان بخش‌ها و زیربخش‌ها
\titlespacing*{\section}{0pt}{1cm}{0.3cm}
\titleformat{\section}
  {\fontsize{14}{6}\bfseries}{\thesection}{1em}{}

\titlespacing*{\subsection}{0pt}{.8cm}{0.2cm}
\titleformat{\subsection}
  {\fontsize{12}{6}\bfseries}{\thesubsection}{1em}{}

\titlespacing*{\subsubsection}{0pt}{.6cm}{0.1cm}
\titleformat{\subsubsection}
  {\fontsize{11}{6}\bfseries}{\thesubsubsection}{1em}{}

% فونت لاتین: Times New Roman، و در صورت نبودن، جایگزین‌های مشابه
\newcommand{\setlatinfontsafe}{%
  \IfFontExistsTF{Times New Roman}{\setlatintextfont{Times New Roman}}{%
  \IfFontExistsTF{TeX Gyre Termes}{\setlatintextfont{TeX Gyre Termes}}{%
  \setlatintextfont{Liberation Serif}}}}
